%% ============================================================================
%% Read, Not Reasoned? Bản trình bày kết quả nghiên cứu cho giảng viên hướng dẫn.
%% Học viên: Nguyễn Nam Hán | GVHD: PGS.TS. Nguyễn Thị Kim Anh
%% Biên dịch: pdfLaTeX hoặc tectonic, chạy 2 lần.
%% Nếu có beamerthemehust.sty cùng thư mục thì dùng theme HUST, nếu không dùng theme đơn giản.
%% Thuật ngữ chuyên ngành giữ tiếng Anh (masking, intervention, ground truth, composition, ...).
%% ============================================================================
\documentclass[aspectratio=169,t,10pt]{beamer}
\usepackage[utf8]{inputenc}
\usepackage[vietnamese]{babel}
\usepackage{amsmath,amssymb}
\usepackage{booktabs}
\usepackage{array}
\usepackage{tabularx}
\usepackage{tikz}
\usetikzlibrary{positioning,arrows.meta}
\usepackage{pgfplots}
\pgfplotsset{compat=1.17,
  boxleg/.style={legend image code/.code={\draw[##1,draw=none] (0cm,-0.07cm) rectangle (0.22cm,0.07cm);}},
  every axis legend/.append style={/tikz/every even column/.append style={column sep=8pt}}}

\definecolor{ink}{HTML}{1F2937}
\definecolor{mid}{HTML}{6B7280}
\definecolor{light}{HTML}{D1D5DB}
\definecolor{bar1}{HTML}{9CA3AF}
\definecolor{bar2}{HTML}{4B5563}
\definecolor{bar3}{HTML}{1F3A5F}

\newif\ifhusttheme
\IfFileExists{beamerthemehust.sty}{\hustthemetrue}{\hustthemefalse}
\ifhusttheme
  \usetheme{hust}
\else
  \usetheme{default}
  \setbeamercolor{structure}{fg=ink}
  \setbeamercolor{normal text}{fg=ink}
  \setbeamercolor{frametitle}{fg=ink}
  \setbeamercolor{title}{fg=ink}
  \setbeamercolor{subtitle}{fg=mid}
  \setbeamertemplate{itemize items}[circle]
  \setbeamertemplate{enumerate items}[default]
  \setbeamerfont{frametitle}{size=\large,series=\bfseries}
  \setbeamertemplate{frametitle}{%
    \vspace{3mm}\hspace*{1mm}\usebeamerfont{frametitle}\insertframetitle\par
    \vspace{1mm}{\color{light}\rule{\textwidth}{0.6pt}}\par\vspace{1mm}}
  \setbeamertemplate{footline}{%
    \hfill{\tiny\color{mid}\insertframenumber\,/\,\inserttotalframenumber}\hspace*{4mm}\vspace*{2mm}}
  \setbeamertemplate{title page}{%
    \vspace*{0.8cm}
    {\LARGE\bfseries\inserttitle\par}
    \vspace{0.3cm}
    {\large\color{mid}\insertsubtitle\par}
    \vspace{0.6cm}
    {\small\insertauthor\par}
    \vspace{0.4cm}
    {\small\color{mid}\insertdate\par}}
  \makeatletter\define@key{beamerframe}{Title}[true]{}\makeatother
\fi
\beamertemplatenavigationsymbolsempty

\newcolumntype{Y}{>{\raggedright\arraybackslash}X}
\newcolumntype{P}[1]{>{\raggedright\arraybackslash}p{#1}}
\newcolumntype{Q}[1]{>{\centering\arraybackslash}p{#1}}
\setlength{\tabcolsep}{5pt}

\newcommand{\ketluan}[1]{\par\vspace{5pt}\noindent{\color{light}\rule{\linewidth}{0.5pt}}\par\vspace{2pt}\noindent{\small\textbf{Kết luận.} #1}}
\newcommand{\nguon}[1]{\par\vspace{2pt}{\tiny\color{mid}#1\par}}
\newcommand{\code}[1]{\texttt{#1}}
\newcommand{\ex}[1]{{\ttfamily\footnotesize #1}}
\newcommand{\dn}{\mathrm{do}_{\text{naive}}}
\newcommand{\dt}{\mathrm{do}_{\text{total}}}

\title[Read, Not Reasoned?]{Read, Not Reasoned?}
\subtitle{Phân tích can thiệp nhân quả có nhãn chuẩn đối với chuỗi suy luận nhiều bước trên mô hình ngôn ngữ nhỏ}
\author[Nguyễn Nam Hán]{\texorpdfstring{%
    \begin{tabular}{rl}
        \textbf{Học viên thực hiện:} & Nguyễn Nam Hán \\[0.3em]
        \textbf{Giảng viên hướng dẫn:} & PGS.TS. Nguyễn Thị Kim Anh
    \end{tabular}}{Nguyễn Nam Hán}%
}
\date{Năm 2026}

\begin{document}

%% ============================================================================
\begin{frame}[noframenumbering,Title]
\titlepage
\end{frame}

\begin{frame}{Nội dung trình bày}
\begin{enumerate}
  \setlength{\itemsep}{6pt}
  \item Bối cảnh, vấn đề, mục tiêu và giả thuyết nghiên cứu
  \item Phương pháp đo: intervention lên chuỗi suy luận với ground truth biết trước
  \item Giai đoạn 1. Chẩn đoán: mô hình dùng Chain-of-Thought như thế nào
  \item Giai đoạn 2. Can thiệp không cần huấn luyện
  \item Giai đoạn 3. Can thiệp bằng fine-tuning
  \item Giai đoạn 4. Ứng dụng: giảm hallucination khi trả lời từ tài liệu (RAG)
  \item Tổng hợp, hạn chế và hướng phát triển
\end{enumerate}
\end{frame}

\section{Bối cảnh và mục tiêu}
%% ============================================================================
\begin{frame}{Bối cảnh và vấn đề}
\textbf{Bối cảnh}
\begin{itemize}
  \item Chain-of-Thought (CoT): mô hình viết các bước trung gian $Z = (z_1, \dots, z_k)$ trước khi đưa ra đáp án $Y$.
  \item Một chuỗi đúng không bảo đảm đáp án được tính từ chuỗi đó. Mô hình có thể trả lời theo lối tắt, còn chuỗi chỉ là phần diễn giải đi kèm (vấn đề faithfulness; Lanham et al., 2023; Paul et al., 2024).
  \item Mô hình nhỏ (0.5B đến 8B) được triển khai rộng rãi vì chi phí thấp và chạy được tại chỗ, nhưng ít được phân tích về mặt này.
\end{itemize}
\vspace{4pt}
\textbf{Vấn đề nghiên cứu}
\begin{itemize}
  \item Các phép đo hiện có (masking một bước, resampling) chưa được kiểm chứng trên dữ liệu có ground truth.
  \item Chưa tách được hai năng lực: đọc thông tin có sẵn trong chuỗi và tự thực hiện composition để tạo thông tin mới.
  \item Chưa rõ nên can thiệp ở giai đoạn decoding hay giai đoạn huấn luyện.
\end{itemize}
\end{frame}

\begin{frame}{Mục tiêu và đóng góp}
\textbf{Mục tiêu}
\begin{itemize}
  \item Xây dựng phép đo mức độ mô hình dùng chuỗi suy luận, kiểm chứng bằng các intervention có ground truth.
  \item Dùng phép đo đó để chẩn đoán mô hình nhỏ, sau đó tìm vị trí can thiệp hiệu quả.
\end{itemize}
\vspace{4pt}
\textbf{Đóng góp}
\begin{enumerate}
  \item \textbf{Phương pháp đo.} Chỉ ra hai phép đo phổ biến cho kết quả lệch, kèm các control và cách hiệu chỉnh.
  \item \textbf{Chẩn đoán.} Mô hình nhỏ đọc đáp án có sẵn trong chuỗi nhưng không dùng chuỗi để thực hiện composition.
  \item \textbf{Can thiệp.} Can thiệp ở bước decoding cuối không hiệu quả. Hai cách hiệu quả: chuyển composition sang rule table, và fine-tuning trên chuỗi.
  \item \textbf{Ứng dụng.} Dữ liệu huấn luyện có tài liệu thiếu bằng chứng và tài liệu mâu thuẫn giúp mô hình biết abstain và bám theo tài liệu.
\end{enumerate}
\end{frame}

\begin{frame}{Giả thuyết nghiên cứu}
\small
\renewcommand{\arraystretch}{1.35}
\begin{tabularx}{\textwidth}{@{}P{0.8cm}YP{2.4cm}@{}}
\toprule
\textbf{Mã} & \textbf{Giả thuyết} & \textbf{Kiểm chứng ở}\\
\midrule
H1 & Các phép đo tầm quan trọng của bước suy luận đang dùng (masking một bước, resampling) cho kết quả lệch khi đối chiếu với ground truth. & TN1, TN2\\
H2 & Mô hình nhỏ chưa fine-tune dùng chuỗi suy luận như nơi đọc đáp án, không như nơi tính ra đáp án. & TN3, TN4\\
H3 & Điểm nghẽn nằm ở composition, nên can thiệp ở bước decoding cuối không hiệu quả còn can thiệp vào composition thì hiệu quả. & TN5 đến TN8\\
H4 & Cùng nguyên tắc intervention có ground truth áp dụng được cho bài toán thực tế: giảm hallucination khi trả lời từ tài liệu. & TN9\\
\bottomrule
\end{tabularx}
\end{frame}

\begin{frame}{Hai loại bài toán suy luận nhiều bước}
\small
\renewcommand{\arraystretch}{1.35}
\begin{tabularx}{\textwidth}{@{}P{2.7cm}YY@{}}
\toprule
 & \textbf{Bridge (MuSiQue)} & \textbf{Composition (CLUTRR)}\\
\midrule
Câu hỏi & Vợ của đạo diễn phim X là ai & Anna có con trai Ben, Ben có chị gái Clara. Clara là gì của Anna\\
Chuỗi suy luận & Hop 1: đạo diễn phim X là Y\newline Hop 2: vợ của Y là Z & Hop 1: con trai của Anna là Ben\newline Hop 2: chị gái của Ben là Clara\\
Đáp án & Z & con gái (daughter)\\
Đáp án có trong chuỗi & Có, ở hop cuối & Không, phải tính son $\circ$ sister = daughter\\
Năng lực cần có & Đọc và chép lại & Composition quan hệ\\
\bottomrule
\end{tabularx}
\vspace{6pt}
\begin{itemize}
  \item Nếu mô hình chỉ đọc, nó sẽ phụ thuộc mạnh vào chuỗi ở MuSiQue và gần như không phụ thuộc ở CLUTRR. So sánh hai loại bài toán cho phép tách việc đọc khỏi việc suy luận.
\end{itemize}
\end{frame}

\section{Phương pháp đo}
%% ============================================================================
\begin{frame}[shrink]{Phương pháp đo: intervention lên chuỗi suy luận}
\small
\textbf{Cách đo.} Không cho mô hình sinh lại. Đưa vào mô hình một prompt cố định gồm ngữ cảnh $X$, chuỗi $Z$ (đã sửa hoặc chưa sửa) và dòng \ex{Answer:}, rồi đọc xác suất mô hình gán cho đáp án đúng.
$$\Delta_i = p(Y^* \mid X, Z) - p\big(Y^* \mid X, \mathrm{do}(z_i)\big)$$
\textbf{Ví dụ trên MuSiQue} (định dạng chuỗi dùng trong thực nghiệm)\par\vspace{2pt}
\footnotesize
\renewcommand{\arraystretch}{1.25}
\begin{tabularx}{\textwidth}{@{}P{3.0cm}Y@{}}
\toprule
\textbf{Phiên bản chuỗi} & \textbf{Nội dung đưa vào sau ngữ cảnh $X$}\\
\midrule
Gốc & \ex{Hop 1: Who directed X? => Y}\newline \ex{Hop 2: Who is the spouse of Y? => Z}\newline \ex{Answer:}\\
$\dn(z_1)$: mask riêng hop 1 & \ex{Hop 1: [MASKED] => [MASKED]}\newline \ex{Hop 2: Who is the spouse of Y? => Z}\\
$\dt(z_1)$: mask hop 1 và các hop dùng kết quả của nó & \ex{Hop 1: [MASKED] => [MASKED]}\newline \ex{Hop 2: Who is the spouse of [MASKED]? => [MASKED]}\\
Không có chuỗi & Chỉ có $X$ và \ex{Answer:}\\
\bottomrule
\end{tabularx}
\small
\begin{itemize}
  \item $\Delta_i$ lớn: mô hình dựa vào $z_i$ khi trả lời. $\Delta_i \approx 0$: mô hình bỏ qua $z_i$.
  \item Ngữ cảnh $X$ luôn giữ nguyên, nên đáp án Z vẫn đọc được từ các đoạn văn.
\end{itemize}
\end{frame}

\begin{frame}[shrink]{Ground truth: các intervention có hiệu ứng thật biết trước}
\small
\textbf{Ý tưởng.} Để biết một phép đo có đúng không, cần những intervention mà hiệu ứng thật đã biết trước nhờ cách xây dựng. Phép đo đúng phải tái hiện được các hiệu ứng này.\par
\vspace{4pt}
\footnotesize
\renewcommand{\arraystretch}{1.3}
\begin{tabularx}{\textwidth}{@{}P{3.0cm}YP{3.1cm}@{}}
\toprule
\textbf{Intervention} & \textbf{Ví dụ trên CLUTRR} & \textbf{Hiệu ứng thật}\\
\midrule
Chèn hop không liên quan & \ex{Hop: Anna likes football.} & Bằng 0\\
Placebo: viết lại hop, giữ nội dung & \ex{Ben is the son of Anna.} thay cho \ex{Anna's son is Ben.} & Bằng 0\\
Mask marker: chèn token che nhưng giữ nội dung & \ex{[MASKED] Anna's son is Ben.} & Chỉ là chi phí của token che\\
Thêm hop ghi sẵn đáp án & \ex{So Clara is Anna's relative: => daughter} & Mô hình biết đáp án nếu đọc\\
Thêm hop ghi đáp án sai & \ex{So Clara is Anna's relative: => mother} & Mô hình suy luận sẽ bỏ qua, mô hình chỉ đọc sẽ chép theo\\
\bottomrule
\end{tabularx}
\small
\begin{itemize}
  \item AUROC: xác suất phép đo chấm hop cần thiết cao hơn hop không liên quan. Giá trị 0.5 tương ứng ngẫu nhiên.
\end{itemize}
\end{frame}

\begin{frame}[shrink]{Các chỉ số đo}
\small
\renewcommand{\arraystretch}{1.3}
\begin{tabularx}{\textwidth}{@{}P{3.3cm}YP{1.9cm}@{}}
\toprule
\textbf{Chỉ số} & \textbf{Ý nghĩa} & \textbf{Dùng ở}\\
\midrule
$\Delta$ & Mức giảm xác suất của đáp án đúng khi can thiệp (đại lượng liên tục). Có thêm bản tính trên log-prob thô để không phụ thuộc tập đáp án ứng viên & TN1 đến TN3, TN8\\
AUROC & Khả năng một phép đo xếp hop cần thiết trên hop không liên quan & TN1\\
Copy-follow rate & Tỷ lệ mô hình chọn đáp án sai được ghi sẵn trong chuỗi & TN3, TN8\\
TE, NDE, TIE & Tác động tổng, trực tiếp và gián tiếp qua chuỗi (mediation analysis) & TN4\\
Accuracy, EM, F1 & Chất lượng đáp án cuối & TN5 đến TN7\\
Hop fidelity & Tỷ lệ hop trong chuỗi tự sinh khớp với đường quan hệ chuẩn & TN7\\
Flip rate & Tỷ lệ đáp án thay đổi khi đổi tên nhân vật, đổi thứ tự câu hoặc thêm distractor & TN7\\
Abstain precision, recall, F1; coverage & Chất lượng của quyết định abstain và tỷ lệ câu được trả lời & TN9\\
\bottomrule
\end{tabularx}
\nguon{Mọi so sánh dùng khoảng tin cậy 95\% bằng paired bootstrap trên cùng tập câu hỏi; so sánh hai phương pháp giải mã dùng thêm kiểm định McNemar.}
\end{frame}

\begin{frame}{Dữ liệu, mô hình và giao thức thực nghiệm}
\small
\renewcommand{\arraystretch}{1.3}
\begin{tabularx}{\textwidth}{@{}P{3.2cm}Y@{}}
\toprule
\textbf{Thành phần} & \textbf{Mô tả}\\
\midrule
CLUTRR & Composition quan hệ gia đình, chuỗi 2 đến 10 hop. Tập kiểm thử gồm câu cùng phân phối, tổ hợp quan hệ mới và chuỗi dài 5 đến 8 hop\\
MuSiQue & Hỏi đáp bridge 2 đến 4 hop trên Wikipedia, mỗi câu kèm 4 hoặc 8 đoạn distractor\\
2WikiMultihopQA & Hỏi đáp nhiều bước trên Wikipedia, dùng để kiểm tra khả năng tổng quát\\
Mô hình & Qwen2.5 (0.5B, 1.5B, 3B, 7B), SmolLM2-1.7B, Qwen3 (1.7B, 4B, 8B)\\
Tài nguyên & 2 GPU T4 16 GB (Kaggle)\\
Giao thức & Giả thuyết và ngưỡng chấp nhận được ghi trước khi chạy (pre-registration). Khoảng tin cậy 95\% bằng paired bootstrap\\
\bottomrule
\end{tabularx}
\end{frame}

\begin{frame}{Các nguyên tắc bảo đảm tính chặt chẽ}
\small
\renewcommand{\arraystretch}{1.35}
\begin{tabularx}{\textwidth}{@{}P{3.4cm}Y@{}}
\toprule
\textbf{Nguyên tắc} & \textbf{Cách thực hiện}\\
\midrule
Ground truth & Mỗi phép đo được kiểm chứng trên intervention có hiệu ứng thật biết trước (slide trước).\\
Pre-registration & Giả thuyết, chỉ số và ngưỡng chấp nhận được ghi trước khi chạy trên tập kiểm thử. Kết quả âm tính được báo cáo đầy đủ.\\
Control & Placebo, mask marker, chuỗi chỉ có hình thức, counterfactual replacement, ablation giữ nguyên đầu vào và chỉ đổi một thành phần.\\
So sánh cùng chi phí & Phương pháp tốn $k$ lượt sinh được so với self-consistency $k$ mẫu có cùng ngân sách token.\\
Kiểm định thống kê & Khoảng tin cậy 95\% bằng paired bootstrap trên cùng tập câu hỏi.\\
Lặp lại & Kết luận chính được kiểm tra trên nhiều cỡ mô hình (0.5B đến 8B), hai họ mô hình và nhiều bộ dữ liệu.\\
\bottomrule
\end{tabularx}
\end{frame}

\begin{frame}{Lộ trình thực nghiệm}
\centering
\begin{tikzpicture}[font=\footnotesize,
  st/.style={draw=mid,rounded corners=3pt,inner sep=5pt,anchor=north west},
  rs/.style={draw=light,fill=black!3,rounded corners=3pt,inner sep=5pt,anchor=north west},
  ar/.style={-{Stealth[length=2.2mm]},thick,mid}]
\node[st] (g1) at (0,0) {\parbox[t][2.35cm][t]{2.95cm}{\raggedright \textbf{Giai đoạn 1}\par\textbf{Chẩn đoán}\par TN1 đến TN4\par\vspace{3pt}Kiểm chứng phép đo, đo mức mô hình dùng chuỗi, phân tích tác động của distractor.}};
\node[st] (g2) at (3.65,0) {\parbox[t][2.35cm][t]{2.95cm}{\raggedright \textbf{Giai đoạn 2}\par\textbf{Can thiệp không huấn luyện}\par TN5, TN6\par\vspace{3pt}Can thiệp ở bước decoding; chuyển composition sang rule table.}};
\node[st] (g3) at (7.3,0) {\parbox[t][2.35cm][t]{2.95cm}{\raggedright \textbf{Giai đoạn 3}\par\textbf{Fine-tuning}\par TN7, TN8\par\vspace{3pt}Fine-tuning trên chuỗi, kiểm chứng trên mô hình mới và dữ liệu thật, đo lại cơ chế.}};
\node[st] (g4) at (10.95,0) {\parbox[t][2.35cm][t]{2.95cm}{\raggedright \textbf{Giai đoạn 4}\par\textbf{Ứng dụng RAG}\par TN9\par\vspace{3pt}Giảm hallucination khi trả lời từ tài liệu.}};
\draw[ar] (g1.east)--(g2.west);
\draw[ar] (g2.east)--(g3.west);
\draw[ar] (g3.east)--(g4.west);
\node[rs] at (0,-2.85) {\parbox[t][1.55cm][t]{2.95cm}{\raggedright \textbf{Phát hiện:} mô hình đọc đáp án có sẵn, không thực hiện composition; distractor tác động qua chuỗi.}};
\node[rs] at (3.65,-2.85) {\parbox[t][1.55cm][t]{2.95cm}{\raggedright \textbf{Phát hiện:} can thiệp ở decoding không hiệu quả; rule table hiệu quả ở miền đóng.}};
\node[rs] at (7.3,-2.85) {\parbox[t][1.55cm][t]{2.95cm}{\raggedright \textbf{Phát hiện:} fine-tuning trên chuỗi hiệu quả nhất và tổng quát hơn.}};
\node[rs] at (10.95,-2.85) {\parbox[t][1.55cm][t]{2.95cm}{\raggedright \textbf{Phát hiện:} mô hình học được cách abstain và bám tài liệu.}};
\end{tikzpicture}
\raggedright
\vspace{4pt}
\begin{itemize}
  \item Phát hiện của mỗi giai đoạn xác định vị trí can thiệp của giai đoạn sau.
\end{itemize}
\end{frame}

\section{Giai đoạn 1. Chẩn đoán}
%% ============================================================================ TN1
\begin{frame}[shrink]{TN1. Masking một hop không phát hiện được hop bridge}
\small
\textbf{Giả thuyết.} Mask riêng một hop giữa ($\dn$, cách đo của Lanham et al., 2023) cho $\Delta \approx 0$, vì hop cuối vẫn ghi sẵn đáp án.\par
\textbf{Thiết kế.} MuSiQue, 150 câu mỗi mô hình (7B: 100 câu). So sánh $\dn$, $\dt$ và bỏ cả chuỗi trên các hop giữa (ví dụ ở slide phương pháp đo).\par
\vspace{4pt}
\footnotesize
\renewcommand{\arraystretch}{1.2}
\begin{tabularx}{\textwidth}{@{}YQ{1.6cm}Q{1.6cm}Q{1.9cm}Q{1.7cm}Q{2.0cm}@{}}
\toprule
\textbf{Mô hình} & $\Delta$ của $\dn$ & $\Delta$ của $\dt$ & Chỉ mask hop cuối & Bỏ cả chuỗi & AUROC của $\dn$\\
\midrule
Qwen2.5-0.5B & $-2.4$ & 56.7 & 67.2 & 52.6 & 0.57\\
Qwen2.5-1.5B & 0.6 & 56.9 & 69.7 & 55.2 & 0.71\\
SmolLM2-1.7B & $-3.9$ & 66.8 & 73.5 & 61.0 & 0.46\\
Qwen2.5-3B & 0.5 & 49.1 & 55.0 & 50.6 & 0.64\\
Qwen2.5-7B & 6.5 & 41.3 & 34.8 & 37.7 & 0.63\\
\bottomrule
\end{tabularx}
\nguon{$\Delta$ tính bằng điểm phần trăm. AUROC phân biệt hop giữa với hop không liên quan.}
\small
\textbf{Kết quả.}
\begin{itemize}
  \item $\dn$ cho $\Delta \approx 0$ và AUROC gần ngẫu nhiên: không phát hiện được hop bridge.
  \item $\dt$, chỉ mask hop cuối và bỏ cả chuỗi cho kết quả xấp xỉ nhau. Trên MuSiQue, $\dt$ thực chất đo việc đáp án có nằm trong chuỗi, không đo riêng vai trò của $z_1$.
\end{itemize}
\ketluan{Masking một hop bị mù với hop bridge. Dù $X$ có sẵn đáp án, mô hình lấy đáp án từ chuỗi suy luận.}
\end{frame}

%% ============================================================================ TN2
\begin{frame}[shrink]{TN2. Resampling với baseline cố định thổi phồng hop không liên quan}
\small
\textbf{Phương pháp resampling} (Bogdan et al., 2025): làm nhiễu hop $z_j$, cho mô hình tự sinh tiếp phần còn lại của chuỗi, rồi đo xác suất đáp án đúng. Baseline thường dùng là chuỗi chuẩn đầy đủ, với $p(Y^*) \approx 1$.\par
\textbf{Giả thuyết.} Mô hình nhỏ tự sinh tiếp thường sai, nên baseline cố định gán cả hop không liên quan là quan trọng.\par
\textbf{Thiết kế.} So sánh baseline cố định với baseline hiệu chỉnh theo vị trí $b(j)$: độ chính xác khi mô hình tự sinh tiếp từ cùng vị trí mà không bị can thiệp.\par
\vspace{4pt}
\begin{columns}[T,onlytextwidth]
\column{0.5\textwidth}
\footnotesize
\renewcommand{\arraystretch}{1.25}
\begin{tabularx}{\linewidth}{@{}YQ{1.0cm}Q{1.0cm}@{}}
\toprule
\textbf{$\Delta$ trên MuSiQue} & \textbf{1.5B} & \textbf{3B}\\
\midrule
Baseline cố định, hop cần thiết & 75.7 & 61.5\\
Baseline cố định, hop không liên quan & 76.3 & 56.3\\
\midrule
Baseline hiệu chỉnh, hop cần thiết & 6.4 & 11.3\\
Baseline hiệu chỉnh, hop không liên quan & $-0.3$ & 4.3\\
\bottomrule
\end{tabularx}
\column{0.46\textwidth}
\small
\textbf{Kết quả.}
\begin{itemize}
  \item Với baseline cố định, hop không liên quan có $\Delta$ ngang hop cần thiết. Độ lệch xấp xỉ $1 - b(j)$, với $b(j)$ từ 22\% đến 44\%.
  \item Baseline hiệu chỉnh đưa hop không liên quan về gần 0.
  \item Hop cần thiết sau hiệu chỉnh cũng nhỏ: mô hình tự tìm lại thông tin từ $X$. Resampling đo khả năng phục hồi, masking đo việc sử dụng.
\end{itemize}
\end{columns}
\ketluan{Resampling chỉ dùng được với baseline hiệu chỉnh theo vị trí, và đo một đại lượng khác với masking.}
\end{frame}

%% ============================================================================ TN3
\begin{frame}[shrink]{TN3. Thiết kế: phân biệt đọc với suy luận trên CLUTRR}
\small
\textbf{Giả thuyết.} Mô hình nhỏ không dùng các hop quan hệ để thực hiện composition, mà chỉ đọc đáp án khi đáp án được viết sẵn trong chuỗi.\par
\textbf{Thiết kế.} 158 câu CLUTRR mỗi mô hình. Chuỗi đưa vào là chuỗi chuẩn (không do mô hình sinh), theo ba điều kiện:\par
\vspace{4pt}
\footnotesize
\renewcommand{\arraystretch}{1.3}
\begin{tabularx}{\textwidth}{@{}P{1.0cm}P{5.6cm}YY@{}}
\toprule
 & \textbf{Chuỗi đưa vào mô hình} & \textbf{Nếu mô hình suy luận} & \textbf{Nếu mô hình chỉ đọc}\\
\midrule
C1 & \ex{Hop 1: Anna's son is Ben. => Ben}\newline\ex{Hop 2: Ben's sister is Clara. => Clara}\newline rồi mask một hop & $\Delta$ lớn, vì mất một quan hệ cần cho composition & $\Delta \approx 0$\\
C2 & C1 thêm hop cuối\newline\ex{So Clara is Anna's relative: => daughter} & $\Delta$ thay đổi ít & $\Delta$ lớn, vì mô hình đọc hop cuối\\
C3 & C1 thêm hop cuối ghi đáp án sai\newline\ex{So Clara is Anna's relative: => mother} & Bỏ qua hop sai & Chép theo hop sai (copy)\\
\bottomrule
\end{tabularx}
\small
\begin{itemize}
  \item Ở C3, nhãn sai được chọn ngẫu nhiên trong các quan hệ khác của CLUTRR. Mỗi cặp người chỉ có đúng một nhãn đúng, nên mọi nhãn khác đều mâu thuẫn với các hop quan hệ.
\end{itemize}
\end{frame}

\begin{frame}[shrink]{TN3. Kết quả C1 và C2: chuỗi chỉ có tác dụng khi chứa sẵn đáp án}
\begin{columns}[T,onlytextwidth]
\column{0.52\textwidth}
\centering
\begin{tikzpicture}
\begin{axis}[ybar, boxleg, width=\linewidth, height=5cm, bar width=8pt, ymin=-5, ymax=85,
  symbolic x coords={0.5B,1.5B,Smol,3B,7B}, xtick=data, x tick label style={font=\scriptsize},
  ylabel={$\Delta$ (điểm \%)}, ylabel style={font=\scriptsize}, y tick label style={font=\scriptsize},
  legend style={at={(0.5,1.03)},anchor=south,legend columns=2,font=\scriptsize,draw=none},
  enlarge x limits=0.12, ymajorgrids, grid style={light}]
\addplot[fill=bar1,draw=none] coordinates {(0.5B,0.1)(1.5B,-1.5)(Smol,1.8)(3B,7.3)(7B,1.3)};
\addplot[fill=bar3,draw=none] coordinates {(0.5B,60.8)(1.5B,76.1)(Smol,69.5)(3B,77.0)(7B,51.9)};
\legend{C1: chỉ có hop quan hệ, C2: có hop ghi đáp án}
\end{axis}
\end{tikzpicture}
\column{0.44\textwidth}
\small
\textbf{Kết quả}
\begin{itemize}
  \item C1: $\Delta$ từ $-1.5$ đến 7.3 điểm. Mask một hop quan hệ gần như không ảnh hưởng.
  \item C2: cùng các hop đó nhưng thêm một dòng ghi đáp án, $\Delta$ tăng lên 52 đến 77 điểm.
  \item Ở 3B, 7.3 [3.8, 11.0] điểm có ý nghĩa thống kê nhưng nhỏ hơn hiệu ứng ở C2 khoảng mười lần.
  \item Khi được đưa chuỗi chuẩn đầy đủ, mô hình chỉ trả lời đúng 15\% (0.5B) đến 33\% (7B).
\end{itemize}
\end{columns}
\ketluan{Các hop quan hệ gần như không tham gia vào câu trả lời. Chuỗi chỉ có tác dụng khi đáp án được viết sẵn trong chuỗi.}
\end{frame}

\begin{frame}[shrink]{TN3. Kết quả C3 (copy test): mô hình chép theo đáp án sai}
\small
\textbf{Lý do cần copy test.} C1 có thể bằng 0 vì mô hình không biết làm composition, không phải vì nó chỉ đọc. Copy test tách hai khả năng này: đo tỷ lệ chép đáp án sai riêng trên những câu mô hình đã trả lời đúng khi chưa có hop sai.\par
\vspace{4pt}
\footnotesize
\renewcommand{\arraystretch}{1.25}
\begin{tabularx}{\textwidth}{@{}YQ{3.0cm}Q{3.4cm}Q{2.6cm}@{}}
\toprule
\textbf{Mô hình} & \textbf{Chép theo, mọi câu} & \textbf{Chép theo, câu tự giải đúng} & \textbf{Số câu tự giải đúng}\\
\midrule
Qwen2.5-0.5B & 66.5\% [58.9, 74.1] & 70.8\% [54.2, 87.5] & 24\\
Qwen2.5-1.5B & 58.2\% [50.6, 66.5] & 63.0\% [44.4, 81.5] & 27\\
SmolLM2-1.7B & 56.3\% [48.1, 64.6] & 57.1\% [42.9, 71.4] & 42\\
Qwen2.5-3B & 46.2\% [38.0, 54.4] & 47.7\% [34.1, 61.4] & 44\\
Qwen2.5-7B & 22.9\% [16.1, 30.5] & 10.3\% [2.6, 20.5] & 39\\
\bottomrule
\end{tabularx}
\nguon{Khoảng tin cậy 95\% bằng bootstrap. Mỗi mô hình 158 câu (7B: 118 câu).}
\small
\textbf{Kết quả.}
\begin{itemize}
  \item Đến 3B, tỷ lệ chép trên câu tự giải đúng ngang với mọi câu: hop ghi sẵn đáp án lấn át kết quả mô hình tự tính được.
  \item Ở 7B, đáp án sai bị bác bỏ khi mô hình tự giải được (10.3\%), nhưng ở C1 các hop quan hệ vẫn không mang hiệu ứng (1.3 điểm). 7B thực hiện composition, nhưng không thông qua chuỗi.
\end{itemize}
\ketluan{Kết quả C1 không do mô hình thiếu năng lực composition. Với mô hình đến 3B, hop ghi sẵn đáp án được ưu tiên hơn kết quả tự tính.}
\end{frame}

\begin{frame}[shrink]{TN3. Có đủ thông tin nhưng không làm composition}
\small
\begin{columns}[T,onlytextwidth]
\column{0.5\textwidth}
\textbf{(a) Chuỗi tự sinh đi đúng đường vẫn không giúp}\par\vspace{2pt}
Tách các câu mà chuỗi CoT do mô hình tự sinh nhắc đủ mọi nhân vật trên đường quan hệ đúng.\par\vspace{3pt}
\footnotesize
\renewcommand{\arraystretch}{1.2}
\begin{tabularx}{\linewidth}{@{}YQ{2.0cm}Q{2.0cm}@{}}
\toprule
\textbf{CoT} & \textbf{Chuỗi đi đủ đường} & \textbf{Chuỗi khác}\\
\midrule
Qwen2.5-1.5B & 17.1\% (n=41) & 18.3\% (n=197)\\
SmolLM2-1.7B & 15.2\% (n=33) & 22.4\% (n=205)\\
\bottomrule
\end{tabularx}
\small
\vspace{6pt}
\textbf{(b) Mô hình trích xuất đúng quan hệ}\par\vspace{2pt}
Khi được yêu cầu liệt kê các quan hệ trong câu chuyện, mô hình tìm được 70\% (1.5B), 71\% (3B) và 87\% (7B) số quan hệ trên đường đúng.
\column{0.46\textwidth}
\textbf{(c) Cùng thông tin, đổi bên thực hiện composition}\par\vspace{2pt}
Giữ nguyên danh sách quan hệ đã trích xuất, chỉ thay bên ghép quan hệ: chính mô hình hoặc rule table. Chuỗi 2 đến 4 hop.\par\vspace{3pt}
\footnotesize
\renewcommand{\arraystretch}{1.2}
\begin{tabularx}{\linewidth}{@{}YQ{1.3cm}Q{1.5cm}@{}}
\toprule
 & \textbf{Mô hình ghép} & \textbf{Rule table ghép}\\
\midrule
Qwen2.5-1.5B & 32.9\% & 46.2\%\\
Qwen2.5-3B & 28.5\% & 54.4\%\\
Qwen2.5-7B & 36.1\% & 71.5\%\\
\bottomrule
\end{tabularx}
\small
\vspace{3pt}
Chênh lệch +13.3, +25.9 và +35.4 điểm, đều có khoảng tin cậy lớn hơn 0.
\end{columns}
\ketluan{Thông tin cần thiết có mặt và được mô hình đọc đúng. Điểm hỏng là bước composition.}
\end{frame}

\begin{frame}[shrink]{Control: loại trừ các giải thích khác cho kết quả masking}
\small
\textbf{Mục đích.} Kiểm tra rằng hiệu ứng đo được đến từ nội dung bị xóa, không đến từ token \ex{[MASKED]}, từ cách viết hay từ hình thức của chuỗi.\par
\vspace{2pt}
\begin{itemize}
  \item \textbf{Hình thức và nội dung:} chuỗi cùng hình dạng nhưng mọi hop bị mask, so với không có chuỗi (hình thức) và so với chuỗi đầy đủ (nội dung).
  \item \textbf{Placebo:} viết lại hop, giữ nguyên nội dung. \textbf{Mask marker:} chèn \ex{[MASKED]} vào cùng các hop như $\dt$ nhưng giữ nội dung.
  \item \textbf{Counterfactual replacement:} thay thực thể bằng một tên khác thay vì mask, không dùng token che.
\end{itemize}
\vspace{2pt}
\footnotesize
\renewcommand{\arraystretch}{1.2}
\begin{tabularx}{\textwidth}{@{}YQ{1.5cm}Q{1.5cm}Q{1.4cm}Q{1.7cm}Q{2.2cm}Q{1.5cm}@{}}
\toprule
\textbf{MuSiQue (điểm)} & \textbf{Hình thức} & \textbf{Nội dung} & \textbf{Placebo} & \textbf{Mask marker} & \textbf{Counterfactual} & $\dt$\\
\midrule
Qwen2.5-0.5B & 4.1 & 48.5 & $-3.1$ & 7.3 & 64.2 & 56.7\\
Qwen2.5-1.5B & 3.3 & 51.9 & 0.2 & 0.0 & 66.5 & 56.9\\
SmolLM2-1.7B & $-5.4$ & 66.4 & 1.1 & 2.3 & 71.6 & 66.8\\
Qwen2.5-3B & 2.3 & 48.3 & 0.3 & $-0.3$ & 61.2 & 49.1\\
Qwen2.5-7B & 2.2 & 35.4 & 0.9 & 5.7 & 45.9 & 41.3\\
\bottomrule
\end{tabularx}
\nguon{Placebo, mask marker, counterfactual và $\dt$ tính trên các hop giữa. CLUTRR có hop ghi đáp án cho cùng kết luận: hình thức từ $-2.2$ đến 0.4, nội dung từ 45.5 đến 75.6.}
\ketluan{Hiệu ứng đến từ nội dung của chuỗi. Hình thức, placebo và token che đều cho hiệu ứng nhỏ; thay thực thể không dùng token che cho cùng kết luận với masking.}
\end{frame}

%% ============================================================================ TN4
\begin{frame}[shrink]{TN4. Distractor gây hại thông qua chuỗi suy luận (mediation analysis)}
\small
\textbf{Giả thuyết.} Đoạn distractor làm giảm độ chính xác chủ yếu vì làm hỏng chuỗi suy luận, không phải vì tác động trực tiếp lên bước trả lời.\par
\textbf{Thiết kế.} MuSiQue, Qwen2.5-3B, 500 câu, thêm 4 đoạn distractor. Ba trạng thái được so sánh:\par
\vspace{3pt}
\centering
\begin{tikzpicture}[font=\footnotesize,
  b/.style={draw=mid,rounded corners=2pt,align=center,inner sep=3pt,minimum height=0.55cm,text width=3.2cm},
  m/.style={draw=mid,dashed,rounded corners=2pt,align=center,inner sep=3pt,minimum height=0.55cm,text width=3.2cm},
  l/.style={align=left,text width=2.2cm,anchor=west},
  ar/.style={-{Stealth[length=2mm]},mid}]
\foreach \y/\lab/\sa/\sb/\ta/\tb in {0/Sạch/b/b/Ngữ cảnh sạch/Chuỗi từ ngữ cảnh sạch, -0.85/Có distractor/m/m/Ngữ cảnh có distractor/Chuỗi sinh lại, -1.7/Giữ chuỗi cũ/m/b/Ngữ cảnh có distractor/Chuỗi từ ngữ cảnh sạch}{
  \node[l] at (0,\y) {\lab};
  \node[\sa] (p) at (4.0,\y) {\ta};
  \node[\sb] (q) at (7.9,\y) {\tb};
  \node[b,text width=1.6cm] (r) at (10.9,\y) {Đáp án};
  \draw[ar] (p)--(q); \draw[ar] (q)--(r);}
\end{tikzpicture}
\raggedright
\vspace{3pt}
\begin{columns}[T,onlytextwidth]
\column{0.42\textwidth}
\footnotesize
\renewcommand{\arraystretch}{1.25}
\begin{tabular}{@{}lcc@{}}
\toprule
\textbf{Chế độ trả lời} & \textbf{NDE} & \textbf{TIE}\\
\midrule
CoT & $-0.6$ & $-5.0$\\
Trả lời trực tiếp & $-3.8$ & $-2.0$\\
\bottomrule
\end{tabular}
\column{0.54\textwidth}
\small
\begin{itemize}
  \item NDE (natural direct effect): thay đổi khi giữ chuỗi cũ, chỉ đổi ngữ cảnh.
  \item TIE (total indirect effect): thay đổi do chuỗi bị sinh lại khác đi.
\end{itemize}
\end{columns}
\ketluan{Với CoT, tác hại của distractor đi qua chuỗi (TIE $-5.0$ điểm, NDE gần 0). Sửa ở bước đọc đáp án sau khi chuỗi đã viết xong không xử lý được nguyên nhân.}
\end{frame}

\section{Giai đoạn 2. Can thiệp không huấn luyện}
%% ============================================================================ TN5
\begin{frame}[shrink]{TN5. Can thiệp ở bước decoding không cải thiện kết quả}
\small
\textbf{Giả thuyết.} Từ TN3 và TN4, can thiệp lúc decoding không sửa được lỗi suy luận. Hai phương pháp được kiểm chứng:
\begin{itemize}
  \item HGCD (đề xuất của luận văn): sau khi chuỗi được sinh, cộng thêm vào điểm của mỗi đáp án ứng viên mức phụ thuộc nhân quả của nó vào các hop.
  \item Lọc ngữ cảnh bằng attention: trước khi sinh chuỗi, chỉ giữ 4 đoạn có trọng số attention từ câu hỏi cao nhất.
\end{itemize}
\vspace{4pt}
\begin{columns}[T,onlytextwidth]
\column{0.48\textwidth}
\footnotesize
\textbf{HGCD, độ chính xác (\%)}\par\vspace{2pt}
\renewcommand{\arraystretch}{1.2}
\begin{tabularx}{\linewidth}{@{}YQ{1.0cm}Q{1.0cm}@{}}
\toprule
 & \textbf{CoT} & \textbf{HGCD}\\
\midrule
CLUTRR, Qwen2.5-1.5B & 18.1 & 18.1\\
CLUTRR, SmolLM2-1.7B & 21.4 & 21.0\\
MuSiQue, Qwen2.5-1.5B & 20.8 & 22.5\\
MuSiQue, SmolLM2-1.7B & 15.0 & 15.8\\
\bottomrule
\end{tabularx}
\column{0.48\textwidth}
\footnotesize
\textbf{Lọc ngữ cảnh, MuSiQue, Qwen2.5-3B (EM \%)}\par\vspace{2pt}
\renewcommand{\arraystretch}{1.2}
\begin{tabularx}{\linewidth}{@{}YQ{1.4cm}Q{1.4cm}@{}}
\toprule
 & \textbf{4 distractor} & \textbf{8 distractor}\\
\midrule
CoT & 38.8 & 31.6\\
Self-consistency 5 mẫu & 41.4 & 35.6\\
Lọc theo attention & 33.8 & 27.4\\
Chỉ giữ đoạn chuẩn & 48.6 & 45.6\\
\bottomrule
\end{tabularx}
\end{columns}
\small
\begin{itemize}
  \item HGCD không khác CoT ($|\Delta| \le 1.7$ điểm, $p \ge 0.6$), đúng như TN3 dự đoán: không có phụ thuộc để khuếch đại.
  \item Lọc theo attention loại nhầm đoạn bridge thứ hai, vì đoạn này chỉ liên quan tới thực thể trung gian chưa xuất hiện trong câu hỏi.
\end{itemize}
\ketluan{Can thiệp ở bước decoding không hiệu quả. Cần tác động trực tiếp vào điểm nghẽn mà TN3 chỉ ra là composition.}
\end{frame}

%% ============================================================================ TN6
\begin{frame}[shrink]{TN6. Chuyển composition sang rule table}
\small
\textbf{Giả thuyết.} Mô hình trích xuất quan hệ tốt nhưng làm composition kém (TN3), nên tách hai việc: mô hình trích xuất, rule table thực hiện composition.\par
\vspace{4pt}
\centering
\begin{tikzpicture}[font=\footnotesize,node distance=0.35cm,
  b/.style={draw=mid,rounded corners=2pt,align=center,minimum height=1.1cm,text width=2.35cm,inner sep=3pt},
  ar/.style={-{Stealth[length=2mm]},mid}]
\node[b] (s) {Câu chuyện};
\node[b,right=of s] (e) {Mô hình trích xuất quan hệ và giới tính};
\node[b,right=of e] (g) {Đồ thị quan hệ giữa các nhân vật};
\node[b,right=of g] (p) {Tìm đường đi giữa hai người được hỏi};
\node[b,right=of p] (r) {Rule table ghép quan hệ trên đường đi};
\draw[ar] (s)--(e); \draw[ar] (e)--(g); \draw[ar] (g)--(p); \draw[ar] (p)--(r);
\end{tikzpicture}
\raggedright
\vspace{4pt}
\textbf{Ví dụ.} Mô hình trích ra \ex{Anna's son is Ben}, \ex{Ben's sister is Clara}. Đường đi Anna $\to$ Ben $\to$ Clara. Rule table: son $\circ$ sister = daughter.\par
\vspace{3pt}
\begin{itemize}
  \item Rule table gồm 126 quy tắc dạng $r_1 \circ r_2 = r_3$, quy nạp tự động từ nhãn của tập huấn luyện (đúng 95\% trên chuỗi chuẩn).
  \item Mô hình trích xuất 5 lần và bỏ phiếu. Đối thủ là self-consistency 10 mẫu, có chi phí sinh token không nhỏ hơn.
  \item Ablation ở TN3 (c) dùng cùng pipeline: giữ nguyên quan hệ đã trích xuất, chỉ thay bên ghép.
\end{itemize}
\end{frame}

\begin{frame}[shrink]{TN6. Kết quả: rule table tăng mạnh độ chính xác}
\begin{columns}[T,onlytextwidth]
\column{0.55\textwidth}
\centering
\begin{tikzpicture}
\begin{axis}[ybar, boxleg, width=\linewidth, height=5.2cm, bar width=5.5pt, ymin=0, ymax=90,
  symbolic x coords={V 1.5B,V 3B,V 7B,S 1.5B,S 3B,S 7B,L 1.5B,L 3B,L 7B}, xtick=data,
  x tick label style={font=\tiny,rotate=35,anchor=east}, y tick label style={font=\scriptsize},
  ylabel={Độ chính xác (\%)}, ylabel style={font=\scriptsize},
  legend style={at={(0.5,1.03)},anchor=south,legend columns=3,font=\scriptsize,draw=none},
  enlarge x limits=0.07, ymajorgrids, grid style={light}]
\addplot[fill=bar1,draw=none] coordinates {(V 1.5B,25.7)(V 3B,36.3)(V 7B,39.7)(S 1.5B,31.6)(S 3B,30.4)(S 7B,39.2)(L 1.5B,34.5)(L 3B,26.0)(L 7B,26.0)};
\addplot[fill=bar2,draw=none] coordinates {(V 1.5B,28.7)(V 3B,38.0)(V 7B,43.3)(S 1.5B,31.0)(S 3B,32.3)(S 7B,40.5)(L 1.5B,29.5)(L 3B,29.0)(L 7B,28.5)};
\addplot[fill=bar3,draw=none] coordinates {(V 1.5B,60.7)(V 3B,70.0)(V 7B,80.0)(S 1.5B,52.5)(S 3B,59.5)(S 7B,74.1)(L 1.5B,33.5)(L 3B,38.5)(L 7B,52.5)};
\legend{CoT,Self-consistency 10,Rule table}
\end{axis}
\end{tikzpicture}
\nguon{V: cùng phân phối, 2 đến 4 hop. S: tổ hợp quan hệ mới. L: chuỗi dài 5 đến 8 hop.}
\column{0.42\textwidth}
\small
\textbf{Kết quả}
\begin{itemize}
  \item Hơn self-consistency 10 mẫu từ 21 đến 37 điểm trên chuỗi 2 đến 4 hop.
  \item Chuỗi dài: hơn 24 điểm ở 7B, không khác biệt ở 1.5B.
  \item Mô hình 1.5B dùng rule table (60.7\%) vượt mô hình 7B dùng self-consistency (43.3\%).
\end{itemize}
\end{columns}
\ketluan{Rule table hiệu quả, và theo ablation ở TN3 (c) phần tăng đến từ composition chứ không từ bước trích xuất. Cách này chỉ áp dụng được ở miền có tập quan hệ đóng.}
\end{frame}

\section{Giai đoạn 3. Fine-tuning}
%% ============================================================================ TN7
\begin{frame}[shrink]{TN7. Fine-tuning trên chuỗi suy luận: thiết kế}
\small
\textbf{Giả thuyết.} Mô hình không làm được composition do cách prompting chứ không do giới hạn năng lực. Fine-tuning trên chuỗi sẽ dạy được mô hình tự làm composition, cả trên chuỗi dài hơn dữ liệu huấn luyện.\par
\textbf{Thiết kế.}
\begin{itemize}
  \item LoRA (rank 16) trên 3000 câu CLUTRR, chỉ gồm chuỗi 2 đến 4 hop. Chuỗi huấn luyện được dựng tự động từ đường quan hệ có sẵn trong dữ liệu.
  \item Kiểm tra trên câu cùng phân phối, tổ hợp quan hệ mới và chuỗi dài 5 đến 8 hop chưa từng thấy khi huấn luyện.
\end{itemize}
\vspace{3pt}
\footnotesize
\renewcommand{\arraystretch}{1.3}
\begin{tabularx}{\textwidth}{@{}P{2.6cm}Y@{}}
\toprule
\textbf{Dạng dữ liệu} & \textbf{Đầu ra mục tiêu cho cùng một câu chuyện}\\
\midrule
Answer-only & \ex{Answer: Clara is Anna's daughter.}\\
Chain & \ex{Hop 1: Anna's son is Ben. => Ben}\newline\ex{Hop 2: Ben's sister is Clara. => Clara}\newline\ex{Answer: Clara is Anna's daughter.}\\
\bottomrule
\end{tabularx}
\small
\begin{itemize}
  \item Hai dạng dùng cùng câu chuyện và cùng số bước huấn luyện, chỉ khác có hay không có các hop.
\end{itemize}
\end{frame}

\begin{frame}[shrink]{TN7. Kết quả: fine-tuning trên chuỗi hiệu quả nhất}
\small
\footnotesize
\renewcommand{\arraystretch}{1.2}
\begin{tabularx}{\textwidth}{@{}YQ{2.2cm}Q{2.2cm}Q{2.2cm}@{}}
\toprule
\textbf{Qwen2.5-3B, độ chính xác (\%)} & \textbf{Cùng phân phối} & \textbf{Tổ hợp mới} & \textbf{Chuỗi dài}\\
\midrule
CoT, không fine-tune & 36.3 & 30.4 & 26.0\\
Rule table (TN6) & 70.0 & 59.5 & 38.5\\
Fine-tune answer-only & 90.3 & 55.1 & 41.0\\
Fine-tune chain & \textbf{94.7} & \textbf{77.8} & \textbf{59.0}\\
\bottomrule
\end{tabularx}
\vspace{4pt}
\begin{columns}[T,onlytextwidth]
\column{0.5\textwidth}
\centering
\begin{tikzpicture}
\begin{axis}[ybar, boxleg, width=\linewidth, height=4.4cm, bar width=6.5pt, ymin=0, ymax=70,
  symbolic x coords={Qwen2.5-1.5B,Qwen2.5-3B,Qwen2.5-7B,Qwen3-4B,Qwen3-8B}, xtick=data,
  x tick label style={font=\tiny,rotate=25,anchor=east}, y tick label style={font=\scriptsize},
  ylabel={Chuỗi dài (\%)}, ylabel style={font=\scriptsize},
  legend style={at={(0.5,1.03)},anchor=south,legend columns=2,font=\scriptsize,draw=none},
  enlarge x limits=0.12, ymajorgrids, grid style={light}]
\addplot[fill=bar1,draw=none] coordinates {(Qwen2.5-1.5B,38.0)(Qwen2.5-3B,41.0)(Qwen2.5-7B,47.5)(Qwen3-4B,40.0)(Qwen3-8B,44.0)};
\addplot[fill=bar3,draw=none] coordinates {(Qwen2.5-1.5B,47.0)(Qwen2.5-3B,59.0)(Qwen2.5-7B,58.5)(Qwen3-4B,61.0)(Qwen3-8B,58.5)};
\legend{Answer-only,Chain}
\end{axis}
\end{tikzpicture}
\column{0.46\textwidth}
\small
\textbf{Kết quả}
\begin{itemize}
  \item Fine-tune chain đưa 3B từ 36.3\% lên 94.7\%, vượt rule table trên chuỗi dài.
  \item Chain hơn answer-only từ 9 đến 21 điểm trên chuỗi dài, lặp lại ở năm mô hình thuộc hai họ.
  \item Trên câu cùng phân phối, answer-only gần ngang chain: khác biệt nằm ở khả năng tổng quát.
  \item Chuỗi do mô hình 3B tự sinh có hop fidelity 0.95 (cùng phân phối) và 0.77 (chuỗi dài). Flip rate: 12\% khi đổi tên, 21\% khi đổi thứ tự câu, 34\% khi thêm distractor.
\end{itemize}
\end{columns}
\ketluan{Mô hình làm được composition sau fine-tuning trên chuỗi, kể cả trên chuỗi dài hơn dữ liệu huấn luyện.}
\end{frame}

\begin{frame}[shrink]{TN7. Kiểm chứng trên mô hình mới và trên dữ liệu thật}
\small
\begin{columns}[T,onlytextwidth]
\column{0.5\textwidth}
\textbf{Qwen3 với chế độ thinking (CLUTRR, \%)}\par\vspace{3pt}
\footnotesize
\renewcommand{\arraystretch}{1.25}
\begin{tabularx}{\linewidth}{@{}YQ{1.2cm}Q{1.2cm}@{}}
\toprule
 & \textbf{Tổ hợp mới} & \textbf{Chuỗi dài}\\
\midrule
Qwen3-4B, few-shot & 38.0 & 27.5\\
Qwen3-4B, thinking & 59.0 & 22.0\\
Qwen3-4B, fine-tune chain & 77.8 & 61.0\\
\midrule
Qwen3-8B, few-shot & 37.3 & 27.5\\
Qwen3-8B, thinking & 60.0 & 31.7\\
Qwen3-8B, fine-tune chain & 83.5 & 58.5\\
\bottomrule
\end{tabularx}
\nguon{Thinking đo trên 100 câu (4B) và 60 câu (8B), giới hạn 1500 token; lượt chưa xong tính là sai.}
\column{0.46\textwidth}
\textbf{2WikiMultihopQA (Wikipedia, Qwen2.5-3B, \%)}\par\vspace{3pt}
\footnotesize
\renewcommand{\arraystretch}{1.25}
\begin{tabularx}{\linewidth}{@{}YQ{1.0cm}Q{1.0cm}@{}}
\toprule
\textbf{Cấu hình} & \textbf{EM} & \textbf{F1}\\
\midrule
Few-shot & 28.8 & 41.9\\
Fine-tune answer-only & 50.0 & 60.4\\
Fine-tune chain & 56.0 & 66.0\\
\midrule
\multicolumn{3}{@{}l}{Huấn luyện trên câu compositional,}\\
\multicolumn{3}{@{}l}{kiểm tra trên câu inference:}\\
Fine-tune answer-only & 22.6 & 45.5\\
Fine-tune chain & 40.6 & 58.9\\
\bottomrule
\end{tabularx}
\end{columns}
\ketluan{Chế độ thinking của mô hình mới chưa thay thế được fine-tuning trên chuỗi. Trên dữ liệu thật, chuỗi giúp tổng quát sang loại câu hỏi chưa thấy khi huấn luyện (+18 điểm).}
\end{frame}

%% ============================================================================ TN8
\begin{frame}[shrink]{TN8. Sau fine-tuning, câu trả lời phụ thuộc vào chuỗi}
\small
\textbf{Thiết kế.} Đo lại TN3 trên mô hình đã fine-tune (120 câu CLUTRR). Trên CLUTRR đáp án không nằm trong chuỗi, nên masking không xóa đáp án như ở TN1.\par
\vspace{4pt}
\footnotesize
\renewcommand{\arraystretch}{1.25}
\begin{tabularx}{\textwidth}{@{}YQ{2.1cm}Q{2.1cm}Q{2.4cm}Q{2.4cm}@{}}
\toprule
 & \multicolumn{2}{c}{\textbf{$\Delta$ khi mask hop giữa (điểm)}} & \multicolumn{2}{c}{\textbf{Chép hop ghi đáp án sai}}\\
\textbf{Mô hình} & \textbf{Trước} & \textbf{Sau} & \textbf{Trước} & \textbf{Sau}\\
\midrule
Qwen2.5-1.5B & $-1.5$ & 36 đến 42 & 63.0\% & 13 đến 22\%\\
Qwen2.5-3B & 7.3 & 41 đến 44 & 47.7\% & 26 đến 32\%\\
Qwen2.5-7B & 1.3 & 39 & 10.3\% & 47.5\%\\
\bottomrule
\end{tabularx}
\nguon{Khoảng giá trị là kết quả của hai seed. Trên chuỗi tuyến tính, $\dt$ cắt mọi liên kết phía sau, nên đo việc dùng chuỗi chứ không cô lập riêng một hop.}
\small
\textbf{Kết quả}
\begin{itemize}
  \item $\Delta$ tăng từ mức nhiễu lên khoảng 40 điểm: sau fine-tuning, câu trả lời phụ thuộc vào chuỗi.
  \item Mô hình vẫn tin hop ghi đáp án ở cuối. Với 7B, tỷ lệ chép đáp án sai tăng sau fine-tuning.
  \item DPO với cặp counterfactual giảm tỷ lệ chép xuống 4.5\% đến 7.2\% nhưng làm giảm 5 đến 13 điểm trên chuỗi dài. GRPO không cải thiện chuỗi dài.
\end{itemize}
\ketluan{Fine-tuning thay đổi cơ chế: mô hình chuyển sang dựa vào chuỗi. DPO và GRPO chỉ cho đánh đổi, không cho cải thiện thêm.}
\end{frame}

\section{Giai đoạn 4. Ứng dụng}
%% ============================================================================ TN9
\begin{frame}[shrink]{TN9. Giảm hallucination khi trả lời từ tài liệu: thiết kế}
\small
\textbf{Giả thuyết.} Thêm vào dữ liệu fine-tuning các tài liệu thiếu bằng chứng và tài liệu mâu thuẫn sẽ dạy mô hình abstain khi tài liệu không đủ và bám theo tài liệu khi tài liệu khác với kiến thức đã nhớ.\par
\vspace{4pt}
\footnotesize
\renewcommand{\arraystretch}{1.3}
\begin{tabularx}{\textwidth}{@{}P{2.4cm}YP{3.3cm}@{}}
\toprule
\textbf{Điều kiện} & \textbf{Intervention lên tài liệu (MuSiQue)} & \textbf{Hành vi đúng}\\
\midrule
Answerable & Giữ nguyên: đoạn bằng chứng và 4 đoạn distractor & Trả lời đúng\\
Unanswerable & Bỏ một đoạn bằng chứng; đáp án không còn xuất hiện trong tài liệu & \ex{Answer: unanswerable}\\
Conflict & Thay đáp án ở mọi chỗ trong tài liệu bằng một thực thể cùng loại & Trả lời theo tài liệu\\
\bottomrule
\end{tabularx}
\small
\vspace{4pt}
\textbf{Ví dụ conflict.} Tài liệu ghi vợ của Y là W thay cho Z. Mô hình bám tài liệu phải trả lời W; trả lời Z nghĩa là dùng kiến thức đã nhớ.\par
\vspace{3pt}
\begin{itemize}
  \item Ba cách fine-tuning: chain thông thường; thêm 30\% mẫu unanswerable; thêm cả mẫu unanswerable và conflict.
  \item Qwen2.5-3B (2 seed) và Qwen3-4B, 300 câu mỗi điều kiện. Ngưỡng chấp nhận đặt trước: độ chính xác trên câu answerable không giảm quá 5 điểm.
\end{itemize}
\end{frame}

\begin{frame}[shrink]{TN9. Kết quả: mô hình học được cách abstain và bám tài liệu}
\footnotesize
\renewcommand{\arraystretch}{1.22}
\setlength{\tabcolsep}{4pt}
\begin{tabularx}{\textwidth}{@{}YQ{1.0cm}Q{1.1cm}Q{1.1cm}Q{1.1cm}Q{1.0cm}Q{1.1cm}Q{1.1cm}Q{1.1cm}@{}}
\toprule
 & \multicolumn{4}{c}{\textbf{Qwen2.5-3B (seed 0)}} & \multicolumn{4}{c}{\textbf{Qwen3-4B}}\\
\cmidrule(lr){2-5}\cmidrule(l){6-9}
\textbf{Fine-tuning} & \textbf{Đúng} & \textbf{Abstain nhầm} & \textbf{Abstain đúng} & \textbf{Theo tài liệu} & \textbf{Đúng} & \textbf{Abstain nhầm} & \textbf{Abstain đúng} & \textbf{Theo tài liệu}\\
\midrule
Không fine-tune, có chỉ dẫn abstain & 29.3 & 17.7 & 35.7 & 11.7 & 31.7 & 35.7 & 73.3 & 13.3\\
Chain thông thường & 52.0 & 0.0 & 0.0 & 22.3 & 60.7 & 0.0 & 0.0 & 30.3\\
+ unanswerable & 44.3 & 19.3 & 61.7 & 21.0 & 56.7 & 12.3 & 70.7 & 32.0\\
+ unanswerable + conflict & 46.7 & 11.7 & 55.7 & 42.7 & 57.3 & 10.7 & 70.7 & 52.7\\
\bottomrule
\end{tabularx}
\nguon{Đúng và abstain nhầm: trên câu answerable. Abstain đúng: trên câu unanswerable. Theo tài liệu: trên câu conflict. Đơn vị \%. Seed 1 của Qwen2.5-3B cho kết quả cùng chiều.}
\small
\vspace{3pt}
\begin{columns}[T,onlytextwidth]
\column{0.48\textwidth}
\textbf{Đạt}
\begin{itemize}
  \item Abstain khi thiếu bằng chứng tăng từ 0\% lên 56\% đến 71\%, lặp lại ở cả 3 lần chạy.
  \item Trả lời theo tài liệu khi conflict tăng 21 đến 24 điểm, lặp lại ở cả 3 lần chạy.
\end{itemize}
\column{0.48\textwidth}
\textbf{Chưa đạt}
\begin{itemize}
  \item Độ chính xác trên câu answerable giảm 4 đến 8 điểm, vượt ngưỡng 5 điểm đặt trước.
  \item DPO và GRPO chỉ dời điểm cân bằng giữa trả lời và abstain, không cải thiện thật.
\end{itemize}
\end{columns}
\ketluan{Hai loại dữ liệu đơn giản dạy được mô hình abstain và bám tài liệu. Nguồn lỗi còn lại chủ yếu là độ chính xác gốc (44\% đến 61\%) trên câu answerable.}
\end{frame}

\begin{frame}[shrink]{TN9. Chỉ số abstention: đánh đổi giữa coverage và độ chính xác}
\small
\textbf{Cách tính.} Lớp dương là câu unanswerable. P, R, F1: precision, recall, F1 của quyết định abstain. Cov.: tỷ lệ câu answerable được trả lời. Ans. P: độ chính xác trên các câu answerable đã trả lời.\par
\vspace{4pt}
\footnotesize
\renewcommand{\arraystretch}{1.2}
\setlength{\tabcolsep}{4pt}
\begin{tabularx}{\textwidth}{@{}YQ{0.6cm}Q{0.6cm}Q{0.6cm}Q{0.95cm}Q{0.95cm}Q{0.6cm}Q{0.6cm}Q{0.6cm}Q{0.95cm}Q{0.95cm}@{}}
\toprule
 & \multicolumn{5}{c}{\textbf{Qwen2.5-3B (seed 0)}} & \multicolumn{5}{c}{\textbf{Qwen3-4B}}\\
\cmidrule(lr){2-6}\cmidrule(l){7-11}
\textbf{Fine-tuning} & \textbf{P} & \textbf{R} & \textbf{F1} & \textbf{Cov.} & \textbf{Ans. P} & \textbf{P} & \textbf{R} & \textbf{F1} & \textbf{Cov.} & \textbf{Ans. P}\\
\midrule
Chain & -- & 0.00 & -- & 1.00 & 0.52 & -- & 0.00 & -- & 1.00 & 0.61\\
+ unans. & 0.76 & 0.62 & 0.68 & 0.81 & 0.55 & 0.85 & 0.71 & 0.77 & 0.88 & 0.65\\
+ unans. + conflict & 0.83 & 0.56 & 0.67 & 0.88 & 0.53 & 0.87 & 0.71 & 0.78 & 0.89 & 0.64\\
\quad + DPO & 0.89 & 0.25 & 0.39 & 0.97 & 0.53 & 0.92 & 0.52 & 0.67 & 0.95 & 0.62\\
\quad + GRPO & 0.73 & 0.81 & 0.77 & 0.70 & 0.58 & 0.73 & 0.90 & 0.81 & 0.68 & 0.72\\
\bottomrule
\end{tabularx}
\nguon{Nguồn: \code{results/hall/abstain\_metrics.csv}. Seed 1 của Qwen2.5-3B cho kết quả cùng chiều.}
\small
\begin{itemize}
  \item Abstain có precision cao (0.76 đến 0.87): khi mô hình abstain, phần lớn đúng là câu thiếu bằng chứng.
  \item DPO giữ coverage cao nhưng recall thấp, tức bỏ sót nhiều câu unanswerable. GRPO tăng recall và answer precision nhưng chỉ trả lời khoảng 70\% câu answerable.
\end{itemize}
\ketluan{Fine-tuning với unanswerable và conflict cho điểm cân bằng tốt nhất. DPO và GRPO dời điểm cân bằng theo hai hướng ngược nhau, không cải thiện đồng thời cả coverage và precision.}
\end{frame}

\section{Tổng hợp và kết luận}
%% ============================================================================
\begin{frame}[shrink]{Tổng hợp kết quả}
\footnotesize
\renewcommand{\arraystretch}{1.18}
\begin{tabularx}{\textwidth}{@{}P{0.9cm}YP{4.8cm}P{1.6cm}@{}}
\toprule
\textbf{TN} & \textbf{Giả thuyết} & \textbf{Kết quả chính} & \textbf{Đánh giá}\\
\midrule
TN1 & Masking một hop không thấy hop bridge & $\Delta \approx 0$, AUROC 0.46 đến 0.71 & Chấp nhận\\
TN2 & Resampling với baseline cố định thổi phồng hop không liên quan & 56 đến 76 điểm, về $\le 4$ sau hiệu chỉnh & Chấp nhận\\
TN3 & Mô hình nhỏ đọc chứ không suy luận & C1 gần 0; C2 từ 52 đến 77 điểm; chép đáp án sai 48\% đến 71\% (đến 3B); control đều nhỏ & Chấp nhận\\
TN4 & Distractor gây hại qua chuỗi & NDE $-0.6$, TIE $-5.0$ & Chấp nhận\\
TN5 & Can thiệp ở decoding cải thiện kết quả & Không khác CoT & Bác bỏ\\
TN6 & Rule table cải thiện kết quả & Hơn self-consistency 21 đến 37 điểm & Chấp nhận\\
TN7 & Fine-tuning trên chuỗi dạy được composition & 3B: 36.3\% lên 94.7\%; chain hơn answer-only 9 đến 21 điểm & Chấp nhận\\
TN8 & Sau fine-tuning mô hình dựa vào chuỗi & $\Delta$ khoảng 40 điểm & Chấp nhận\\
TN9 & Dữ liệu unanswerable và conflict giảm hallucination & Abstain 56\% đến 71\%; theo tài liệu tăng 21 đến 24 điểm; độ chính xác giảm 4 đến 8 điểm & Chấp nhận một phần\\
\bottomrule
\end{tabularx}
\end{frame}

\begin{frame}{Hạn chế}
\begin{itemize}
  \setlength{\itemsep}{5pt}
  \item Dữ liệu chủ yếu là CLUTRR (tổng hợp) và MuSiQue. Rule table ở TN6 chỉ dùng được ở miền có tập quan hệ đóng.
  \item Mô hình từ 0.5B đến 8B. Chưa kiểm chứng trên mô hình rất lớn như GPT-5 hay Claude Opus.
  \item Nhiều kết quả chỉ có một seed, và số câu tự giải đúng trong copy test nhỏ (24 đến 44 câu). Các kết luận chính đã được lặp lại trên nhiều mô hình.
  \item Trên MuSiQue, $\dt$ đo việc đáp án có trong chuỗi, không đo riêng vai trò từng hop. Kết luận về suy luận dựa trên CLUTRR.
  \item Nhãn unanswerable ở TN9 được dựng bằng cách bỏ đoạn văn, mô hình vẫn có thể trả lời từ kiến thức đã nhớ.
  \item Chưa đo calibration của xác suất mô hình, và chưa hiệu chỉnh cho nhiều phép so sánh (Holm, Bonferroni). Các kết luận chính có khoảng tin cậy cách xa 0.
\end{itemize}
\end{frame}

\begin{frame}{Kết luận và hướng phát triển}
\textbf{Kết luận}
\begin{enumerate}
  \item Masking một hop và resampling với baseline cố định cho kết quả lệch. Cần control và baseline hiệu chỉnh.
  \item Mô hình nhỏ chưa fine-tune đọc đáp án có sẵn trong chuỗi nhưng không dùng chuỗi để làm composition, dù có đủ thông tin.
  \item Can thiệp ở bước decoding không hiệu quả. Rule table hiệu quả ở miền đóng. Fine-tuning trên chuỗi hiệu quả và tổng quát hơn.
  \item Dữ liệu unanswerable và conflict giúp mô hình abstain và bám tài liệu khi trả lời từ tài liệu.
\end{enumerate}
\vspace{4pt}
\textbf{Hướng phát triển}
\begin{itemize}
  \item Nâng độ chính xác gốc khi trả lời từ tài liệu bằng dữ liệu chuỗi lớn và đa dạng hơn.
  \item Kiểm chứng với hệ thống RAG dùng retriever thật.
  \item Áp dụng copy test và các intervention có ground truth cho mô hình lớn qua API.
\end{itemize}
\end{frame}

\begin{frame}[c,noframenumbering]
\centering
{\Large\bfseries Xin trân trọng cảm ơn\par}
\vspace{0.4cm}
{Kính mong nhận được ý kiến đóng góp của thầy cô\par}
\end{frame}

\appendix
\begin{frame}{Phụ lục. Tài liệu tham khảo chính}
\small
\begin{enumerate}
  \setlength{\itemsep}{6pt}
  \item Lanham et al. Measuring Faithfulness in Chain-of-Thought Reasoning. arXiv:2307.13702, 2023.
  \item Dziri et al. Faith and Fate: Limits of Transformers on Compositionality. NeurIPS 2023.
  \item Paul, West, Bosselut, Faltings. Making Reasoning Matter: Measuring and Improving Faithfulness of Chain-of-Thought Reasoning. Findings of EMNLP 2024.
  \item Bogdan, Macar, Nanda, Conmy. Thought Anchors: Which LLM Reasoning Steps Matter. arXiv:2506.19143, 2025.
  \item Somov, Chaichuk, Seleznyov, Panchenko, Tutubalina. Breaking the Chain: A Causal Analysis of LLM Faithfulness to Intermediate Structures. arXiv:2603.16475, 2026.
\end{enumerate}
\end{frame}

\begin{frame}{Phụ lục. Mã nguồn và kết quả}
\small
\renewcommand{\arraystretch}{1.25}
\begin{tabularx}{\textwidth}{@{}P{5.2cm}Y@{}}
\toprule
\textbf{Nội dung} & \textbf{Đường dẫn}\\
\midrule
Giả thuyết đăng ký trước và kết luận từng vòng & \code{REVIEW.md}\\
TN1 đến TN3, control & \code{results/reliance\_v2/}, \code{tools/check\_r1\_triviality.py}\\
TN4 & \code{results/mediation\_n500/}\\
TN5, TN6 & \code{results/decoding/}, \code{results/causaldec/}\\
TN7, TN8 & \code{results/cftrain/}, \code{results/cfrl/}, \code{results/egen/}\\
TN9 & \code{hgcd/hall.py}, \code{results/hall/}, \code{tools/analyze\_hall\_abstain.py}\\
\bottomrule
\end{tabularx}
\end{frame}

\end{document}
